\section{動的保持延長}

\begin{definition}[追加保持コスト関数]
$n_{\mathrm{ext}}\in\mathbb{Z}_{\ge0}$，バックアップ間隔 $\Delta>0$，
パラメータ $c_n>0,\ L_{\mathrm{fail}}>0,\ \alpha>0,\ t_m>0$，および
実効保持期間
\[
  T(n_{\mathrm{ext}})=(n+n_{\mathrm{ext}})\Delta-\tau_R
\]
とする．実行可能集合を
\[
  \mathcal{N}=\{m\in\mathbb{Z}_{\ge0}:T(m)>t_m\}
\]
で定める．追加保持コストを $c_n n_{\mathrm{ext}}$，
復旧不能ペナルティを
$L_{\mathrm{fail}}(t_m/T(n_{\mathrm{ext}}))^\alpha$ とすれば，
\[
  B(n_{\mathrm{ext}})
  =c_n n_{\mathrm{ext}}
  +L_{\mathrm{fail}}\left(\frac{t_m}{T(n_{\mathrm{ext}})}\right)^\alpha ,
  \qquad n_{\mathrm{ext}}\in\mathcal{N}.
\]
\end{definition}

\begin{proposition}[最適追加保持世代数]
$B$ は連続緩和の実行可能領域
$\{x\ge0:(n+x)\Delta-\tau_R>t_m\}$ 上で狭義凸である．
制約を無視した停留点は
\[
  T^\ast=\left(\frac{\alpha L_{\mathrm{fail}}t_m^\alpha\Delta}{c_n}
  \right)^{1/(\alpha+1)},\qquad
  n_{\mathrm{ext}}^c=\frac{T^\ast+\tau_R}{\Delta}-n .
\]
$\Delta>0$ なので $\mathcal{N}$ は空でなく，$m_0=\min\mathcal{N}$ とし，
$x^\ast=\max(m_0,n_{\mathrm{ext}}^c)$ とおく．整数最適解は
\[
  n_{\mathrm{ext}}^\ast
  \in\arg\min_{m\in\mathcal{C}} B(m),\qquad
  \mathcal{C}=\{\lfloor x^\ast\rfloor,\lceil x^\ast\rceil\}\cap\mathcal{N}
\]
で与えられる．
\end{proposition}

\begin{proof}
$T=(n+n_{\mathrm{ext}})\Delta-\tau_R$ として微分すると
\[
  \frac{\dd B}{\dd n_{\mathrm{ext}}}
  =c_n-L_{\mathrm{fail}}\alpha t_m^\alpha T^{-(\alpha+1)}\Delta .
\]
一階条件から $T=T^\ast$ を得る．また
\[
  \frac{\dd^2 B}{\dd n_{\mathrm{ext}}^2}
  =L_{\mathrm{fail}}\alpha(\alpha+1)t_m^\alpha T^{-(\alpha+2)}\Delta^2>0
\]
なので狭義凸である．従って連続緩和の制約付き最小点は
$x^\ast$ であり，整数格子上の最小点はその左右の近傍整数だけを
評価すればよい．
\end{proof}
